% Part: normal-modal-logic
% Chapter: filtrations
% Section: S5-decidable

\documentclass[../../../include/open-logic-section]{subfiles}

\begin{document}

\olfileid{nml}{fil}{dec}

\olsection{\Log{S5} నిర్ణయించదగినది}
స్కీమాల ద్వారా ఇచ్చిన మోడల్ తర్క వ్యవస్థలు
\emph{నిర్ణయించదగినవి} అని చూపడానికి పరిమిత నమూనా
ధర్మం సులభమైన మార్గాన్ని ఇస్తుంది. అంటే,
\tetoken{ఒక సూత్రం}{!!a{formula}} వ్యవస్థలో
\tetoken{వ్యుత్పాద్యమో}{!!{derivable}} కాదో
నిర్ణయించే గణనీయ ప్రక్రియ ఉంటుంది.

\begin{thm}
  \Log{S5} నిర్ణయించదగినది.
\end{thm}

\begin{proof}
  $!A$ ఇవ్వబడిందనుకుందాం; $!A$లో వచ్చే
  \tetoken{ప్రతిజ్ఞావాక్య చరరాశులు}{!!{propositional variable}s}
  $p_1$, \dots, $p_k$లలో ఉన్నాయనుకుందాం.
  ప్రతి $n$కూ, $n$ లోకాలతో $p_1$, \dots, $p_k$లకు
  విలువలను కేటాయించే సార్వత్రిక నమూనాలు పరిమిత
  సంఖ్యలోనే ఉంటాయి. అందువల్ల \emph{సమాంతరంగా}
  రెండు ప్రక్రియలు నడపవచ్చు: ఒక క్రమబద్ధమైన
  విధానంలో నిరూపణలను ఉత్పత్తి చేసి
  \Log{S5} సిద్ధాంతాలన్నింటినీ జాబితా చేయడం;
  అలాగే $1$, $2$, \dots~లోకాలున్న పరిమిత
  సార్వత్రిక నమూనాలన్నింటినీ జాబితా చేసి,
  వాటిలో ఏదో లోకం వద్ద $!A$ అసత్యమో పరీక్షించడం.
  \sourcecorrection{OLTENMLFILDEC-001}{స్థిర మూలం
    ఇక్కడ ``అన్ని నమూనాలు'' అంటుంది. S5కు సంబంధం
    లేని పరిమిత నమూనాలో ప్రతినమూనా దొరికితే
    నిర్ణయ ప్రక్రియ తప్పుడు ప్రతికూల జవాబు ఇస్తుంది.
    అందువల్ల ముందరి S5 పరిమిత నమూనా నిరూపణకు
    అనుగుణంగా సార్వత్రిక నమూనాలకే శోధనను
    పరిమితం చేశాం.}
  ఈ రెండు సమాంతర ప్రక్రియల్లో ఒకటి చివరకు
  సమాధానం ఇస్తుంది: \olref[com][fra]{thm:generaldet},
  \olref[fmp]{cor:S5fmp} ప్రకారం, $!A$
  \tetoken{వ్యుత్పాద్యం}{!!{derivable}} లేదా అది
  ఒక పరిమిత సార్వత్రిక నమూనాలో అసత్యం.
\end{proof}

పైన చెప్పిన నిరూపణ \Log{S5}కు పనిచేయడానికి కారణం,
సార్వత్రిక నమూనాల వడపోతలు స్వయంగా సార్వత్రికమే
కావడం. స్వావర్తనత్వం, సీరియల్ ధర్మాల విషయంలోనూ
ఇది వర్తిస్తుంది; ఇతర ధర్మాలకు ఇంకా పని అవసరం.

\end{document}
